\section{Introduction}\label{sec:introduction}

\subsection{Physics-informed neural networks and the certification problem}

Physics-informed neural networks (PINNs)~\cite{raissi2019physics} have
emerged as a flexible framework for solving forward and inverse problems
governed by partial differential equations (PDEs). In this approach, a neural
network approximation $\hat{u}$ is trained by minimizing a composite loss
functional that penalizes violations of the governing PDE in the interior
and mismatches of boundary conditions on $\partial\Omega$. The appeal of
PINNs lies in their meshfree nature, the use of automatic differentiation
for exact derivative evaluation, and their applicability across a wide range
of PDE classes~\cite{karniadakis2021physics}.

In most practical implementations, boundary conditions are enforced through
\emph{soft penalty terms} rather than strong constraints. Consequently, the
trained approximation $\hat{u}$ does not, in general, satisfy the boundary
conditions exactly.

Despite their empirical success, a fundamental limitation of PINNs remains:
\emph{the lack of reliable error certification}. In practice, the training
loss is often used as a proxy for solution accuracy. However, this quantity
does not provide control of the approximation error in the energy seminorm.
Indeed, PINNs may exhibit very small training loss while retaining
substantial error in
\[
|u - \hat{u}|_{H^1(\Omega)} = \|\nabla(u - \hat{u})\|_{L^2(\Omega)}.
\]
This discrepancy is not merely a consequence of incomplete optimization,
but a structural feature of the soft-penalty formulation used to enforce
boundary conditions.

Providing a \emph{computable and reliable upper bound} for the approximation
error of PINNs is therefore a central open problem. The present work addresses
this issue by developing an a posteriori error estimator that is fully
computable from a trained network and does not require access to the exact
solution.

\subsection{Classical a posteriori error estimation}

For elliptic PDEs discretized by the finite element method (FEM), a posteriori
error estimation is a well-established theory. Classical residual-based
estimators provide bounds of the form
\begin{equation}\label{eq:classical-aposteriori}
|u - u_h|_{H^1(\Omega)} \leq C_{\mathrm{rel}}\, \eta(u_h),
\end{equation}
where $\eta(u_h)$ is a computable quantity assembled from element residuals
and flux jumps; see Ainsworth and Oden~\cite{ainsworth1997posteriori} and
Verf\"urth~\cite{verfurth2013posteriori}. When the dual norm is evaluated
exactly, the reliability constant can be taken as $C_{\mathrm{rel}} = 1$.

In contrast, no directly comparable framework exists for PINNs. While the
residual is naturally available through automatic differentiation, it is
typically monitored in an empirical $L^2$ sense, and the boundary mismatch is
treated heuristically via penalty terms. Existing theoretical results for
PINNs~\cite{mishra2022estimates, deryck2022generic, shin2020convergence}
provide convergence and generalization estimates, but do not yield
\emph{computable a posteriori bounds} that can be evaluated from a trained
network.

\subsection{The boundary mismatch problem}

A central observation of this work is that residual-based estimators,
when applied directly to PINNs, fail to capture the dominant component of
the error when boundary conditions are enforced weakly.

Let $e := u - \hat{u}$ denote the approximation error. Since $\hat{u}$ does not
satisfy the Dirichlet boundary condition exactly, we generally have
$e \notin H^1_0(\Omega)$. We therefore introduce the decomposition
\[
e = e_0 + w,
\]
where $e_0 \in H^1_0(\Omega)$ and $w$ is the harmonic lifting of the boundary
mismatch.

By the triangle inequality,
\[
|e|_{H^1(\Omega)} \leq |e_0|_{H^1(\Omega)} + \|\nabla w\|_{L^2(\Omega)}.
\]

The second term represents the contribution of boundary mismatch propagated
into the domain. Numerical evidence demonstrates that this term dominates the
total error across all validated benchmarks. In particular, residual-only
estimators severely underestimate the true error, while the inclusion of the
boundary lifting term restores reliability.

\subsection{Contributions}

The contributions of this work are as follows.

\begin{enumerate}[leftmargin=*]

\item \textbf{Residual-based estimator for PINNs.}
We derive a computable a posteriori error estimator of the form
\[
|u - \hat{u}|_{H^1(\Omega)}
\leq
\frac{1}{\alpha}\|R\|_{H^{-1}(\Omega)} + \|\nabla w\|_{L^2(\Omega)},
\]
where
\[
R = f + \nabla \cdot (a \nabla \hat{u})
\]
is the interior residual and $w$ is the harmonic lifting of the boundary mismatch.

\item \textbf{Reliability analysis.}
We prove reliability for constant diffusion coefficients and provide a
computable and practically reliable upper bound for variable-coefficient
problems using a Laplace lifting surrogate.

\item \textbf{Computational realization.}
We provide an efficient implementation based on Discrete Sine Transform (DST)
methods for rectangular domains and finite element discretization for
nontrivial geometries.

\item \textbf{Experimental validation.}
A reproducible benchmark suite demonstrates effectivity indices in the range
$1.05 \le \eta \le 1.31$ across smooth, heterogeneous, and singular problems,
confirming reliability and sharpness of the estimator.

\item \textbf{Failure of residual-only certification.}
We demonstrate that residual-only estimators can underestimate the error by
orders of magnitude, whereas the proposed estimator accurately captures the
true error through the inclusion of boundary mismatch.

\end{enumerate}